% 2008 年广东卷理科第 7 题：正三棱柱截去三个角前后的图 1、图 2，以及四个侧视图选项。
% 来源 PDF：gaokaomath-src/普通高考/2008/2008广东理.pdf
% SHA-256：c395b9608135bcd958aa64c448dfc14a8ba00fbc84cd20c3f1ad316e94d549b3
% 原页：2008广东理.pdf 第 1 页（A3 三栏）第 1 栏，书页码 445。
% 重绘依据：图 1 为正三棱柱 GHI-EFD（H、G 在后，I 在前），A、B、C 分别是 HG、HI、GI 的中点，
%   截角平面为 ABE、ACD、BCF，故新增棱 AB、AC、BC、BE、CD、BF、CF，被遮挡的 AE、AD、ED 用虚线；
%   图 2 为截角后的几何体，只剩 A、B、C、D、E、F 六个顶点，实线为 AB、AC、BC、BE、CD、BF、CF、
%   EF、FD，虚线为 AE、AD、ED。
\begin{tikzpicture}[semithick, font=\small, every node/.style={inner sep=1.2pt}]

  % ---------------- 图 1：正三棱柱 GHI-EFD 与截角线 ----------------
  \begin{scope}[shift={(0,0)}]
    \coordinate (H) at (0,3.45);
    \coordinate (A) at (1.90,3.45);
    \coordinate (G) at (3.80,3.45);
    \coordinate (I) at (1.90,2.15);
    \coordinate (B) at (0.95,2.80);
    \coordinate (C) at (2.85,2.80);
    \coordinate (E) at (0,0.40);
    \coordinate (D) at (3.80,0.40);
    \coordinate (F) at (1.90,-0.85);

    % 可见棱
    \draw (H) -- (G);
    \draw (H) -- (B);
    \draw (G) -- (C);
    \draw (H) -- (E);
    \draw (G) -- (D);
    \draw (B) -- (I) -- (C);
    \draw (I) -- (F);
    \draw (E) -- (F) -- (D);
    % 截角线（可见部分）
    \draw (A) -- (B);
    \draw (A) -- (C);
    \draw (B) -- (C);
    \draw (B) -- (E);
    \draw (C) -- (D);
    \draw (B) -- (F);
    \draw (C) -- (F);
    % 被遮挡的截角线与底边
    \draw[dashed] (A) -- (E);
    \draw[dashed] (A) -- (D);
    \draw[dashed] (E) -- (D);

    \node[above] at (H) {$H$};
    \node[above] at (A) {$A$};
    \node[above] at (G) {$G$};
    \node[below left=1pt] at (B) {$B$};
    \node[below right=1pt] at (C) {$C$};
    \node[right=3pt] at (I) {$I$};
    \node[left] at (E) {$E$};
    \node[right] at (D) {$D$};
    \node[below] at (F) {$F$};
    \node at (1.90,-1.75) {图 1};
  \end{scope}

  % 侧视方向箭头
  \draw[->,>=stealth] (4.45,1.55) -- (5.45,1.55);
  \node[above] at (4.95,1.55) {侧视};

  % ---------------- 图 2：截角后的几何体 ----------------
  \begin{scope}[shift={(6.55,-0.05)}]
    \coordinate (A2) at (1.90,3.35);
    \coordinate (B2) at (0.95,2.62);
    \coordinate (C2) at (2.85,2.62);
    \coordinate (E2) at (0.10,0.30);
    \coordinate (D2) at (3.70,0.30);
    \coordinate (F2) at (1.90,-0.95);

    % 可见棱
    \draw (A2) -- (B2) -- (C2) -- cycle;
    \draw (B2) -- (E2);
    \draw (C2) -- (D2);
    \draw (B2) -- (F2);
    \draw (C2) -- (F2);
    \draw (E2) -- (F2) -- (D2);
    % 被遮挡的棱
    \draw[dashed] (A2) -- (E2);
    \draw[dashed] (A2) -- (D2);
    \draw[dashed] (E2) -- (D2);

    \node[above] at (A2) {$A$};
    \node[below left=1pt] at (B2) {$B$};
    \node[below right=1pt] at (C2) {$C$};
    \node[left] at (E2) {$E$};
    \node[right] at (D2) {$D$};
    \node[below] at (F2) {$F$};
    \node at (1.90,-1.75) {图 2};
  \end{scope}

  % ---------------- 四个侧视图选项 ----------------
  % (A) 左边竖直、顶边半宽、右侧斜边，内部连 E—B
  \begin{scope}[shift={(0,-6.10)}]
    \draw (0,0) -- (0,2.60) -- (1.20,2.60) -- (2.60,0) -- cycle;
    \draw (0,0) -- (1.20,2.60);
    \node[above] at (1.20,2.60) {$B$};
    \node[left] at (0,0) {$E$};
    \node at (1.30,-0.55) {(A)};
  \end{scope}
  % (B) 两侧均为斜边，B 在顶边右端，内部连 E—B
  \begin{scope}[shift={(4.60,-6.10)}]
    \draw (0,0) -- (0.90,2.40) -- (1.55,2.40) -- (2.60,0) -- cycle;
    \draw (0,0) -- (1.55,2.40);
    \node[above] at (1.55,2.40) {$B$};
    \node[left] at (0,0) {$E$};
    \node at (1.30,-0.55) {(B)};
  \end{scope}
  % (C) 右边竖直、左侧斜边，内部连 E—右上角
  \begin{scope}[shift={(0,-9.70)}]
    \draw (0,0) -- (0.90,2.60) -- (2.60,2.60) -- (2.60,0) -- cycle;
    \draw (0,0) -- (2.60,2.60);
    \node[above left] at (0.90,2.60) {$B$};
    \node[left] at (0,0) {$E$};
    \node at (1.30,-0.55) {(C)};
  \end{scope}
  % (D) 两侧均为斜边、整体更窄，B 在顶边右端，内部连 E—B
  \begin{scope}[shift={(4.60,-9.70)}]
    \draw (0,0) -- (0.45,2.40) -- (1.05,2.40) -- (2.60,0) -- cycle;
    \draw (0,0) -- (1.05,2.40);
    \node[above] at (1.05,2.40) {$B$};
    \node[left] at (0,0) {$E$};
    \node at (1.30,-0.55) {(D)};
  \end{scope}
\end{tikzpicture}
